% multi-claim.tex — fixture: paper with multiple claims and no edges.
%
% Purpose: exercise ordinal-pairing in the parser. The Nth claim env in
% source must pair with the Nth RRXIV:claim:N marker in the sidecar.

\documentclass{rrxiv}

\rrxivid{rrxiv-fixture-multi-claim}
\rrxivversion{v1}
\rrxivprotocolversion{0.1.0}
\rrxivlicense{CC-BY-4.0}
\rrxivtopics{example,conformance}

\title{Conformance fixture: multi-claim paper}
\author{rrxiv Project}

\begin{document}
\maketitle

\begin{abstract}
Three claims of three different types (theoretical, empirical,
methodological) so the parser's ordinal-pairing can be exercised
without any inline edges complicating the picture.
\end{abstract}

\section{Theoretical claim}

\begin{claim}[Existence]
\label{claim:existence}
There exists at least one rrxiv paper with three or more semantic
environments of the same kind.
\end{claim}

\section{Empirical claim}

\begin{claim}[Volume]
\label{claim:volume}
The number of rrxiv papers grows monotonically with time once
submissions are accepted.
\end{claim}

\section{Methodological claim}

\begin{claim}[Pairing]
\label{claim:pairing}
The Nth \texttt{claim} environment block in source pairs with the Nth
\texttt{RRXIV:claim:N} marker in the emitted sidecar.
\end{claim}

\end{document}
